\documentclass[10pt,a4paper]{amsart}
\usepackage{amsmath,amssymb,amsthm,graphicx}
% Retained preamble for the Stacks Project chapter {CH} (shared project preamble).
% The chapter file pinned-original-source.tex only carries \input{{preamble}} on its first line,
% so no preamble text lives inside the chapter itself.  The Stacks Project uses one single shared
% preamble for every chapter: the snapshot stages it three times (sources/stacks-chow/preamble.tex,
% sources/stacks-etale-cohomology/preamble.tex and focused-seeds/AGP/sources/stacks-snapshot-20261011/
% preamble.tex) and those copies are identical in size.  The block below is copied mechanically,
% line for line, from sources/stacks-chow/preamble.tex lines 14-18 and 154-217: the Xy-pic setup and
% the native theorem declarations together with the mathematical macros.  Nothing was added, removed
% or redefined; the wrapper class and the amsmath/amssymb/amsthm/graphicx base packages are supplied
% by the assembler.
% For commutative diagrams we use Xy-pic
% Chapter: stacks-cohomology.tex -- Cohomology of Sheaves.
\usepackage[all]{xy}

% We use 2cell for 2-commutative diagrams.
\xyoption{2cell}

% Theorem environments.
%
\theoremstyle{plain}
\newtheorem{theorem}[subsection]{Theorem}
\newtheorem{proposition}[subsection]{Proposition}
\newtheorem{lemma}[subsection]{Lemma}

\theoremstyle{definition}
\newtheorem{definition}[subsection]{Definition}
\newtheorem{example}[subsection]{Example}
\newtheorem{exercise}[subsection]{Exercise}
\newtheorem{situation}[subsection]{Situation}

\theoremstyle{remark}
\newtheorem{remark}[subsection]{Remark}
\newtheorem{remarks}[subsection]{Remarks}

\numberwithin{equation}{subsection}

% Macros
%
\def\lim{\mathop{\mathrm{lim}}\nolimits}
\def\colim{\mathop{\mathrm{colim}}\nolimits}
\def\Spec{\mathop{\mathrm{Spec}}}
\def\Hom{\mathop{\mathrm{Hom}}\nolimits}
\def\Ext{\mathop{\mathrm{Ext}}\nolimits}
\def\SheafHom{\mathop{\mathcal{H}\!\mathit{om}}\nolimits}
\def\SheafExt{\mathop{\mathcal{E}\!\mathit{xt}}\nolimits}
\def\Sch{\mathit{Sch}}
\def\Mor{\mathop{\mathrm{Mor}}\nolimits}
\def\Ob{\mathop{\mathrm{Ob}}\nolimits}
\def\Sh{\mathop{\mathit{Sh}}\nolimits}
\def\NL{\mathop{N\!L}\nolimits}
\def\CH{\mathop{\mathrm{CH}}\nolimits}
\def\proetale{{pro\text{-}\acute{e}tale}}
\def\etale{{\acute{e}tale}}
\def\QCoh{\mathit{QCoh}}
\def\Ker{\mathop{\mathrm{Ker}}}
\def\Im{\mathop{\mathrm{Im}}}
\def\Coker{\mathop{\mathrm{Coker}}}
\def\Coim{\mathop{\mathrm{Coim}}}

% Boxtimes
%
\DeclareMathSymbol{\boxtimes}{\mathbin}{AMSa}{"02}

%
% Macros for moduli stacks/spaces
%
\def\QCohstack{\mathcal{QC}\!\mathit{oh}}
\def\Cohstack{\mathcal{C}\!\mathit{oh}}
\def\Spacesstack{\mathcal{S}\!\mathit{paces}}
\def\Quotfunctor{\mathrm{Quot}}
\def\Hilbfunctor{\mathrm{Hilb}}
\def\Curvesstack{\mathcal{C}\!\mathit{urves}}
\def\Polarizedstack{\mathcal{P}\!\mathit{olarized}}
\def\Complexesstack{\mathcal{C}\!\mathit{omplexes}}
% \Pic is the operator that assigns to X its picard group, usage \Pic(X)
% \Picardstack_{X/B} denotes the Picard stack of X over B
% \Picardfunctor_{X/B} denotes the Picard functor of X over B
\def\Pic{\mathop{\mathrm{Pic}}\nolimits}
\def\Picardstack{\mathcal{P}\!\mathit{ic}}
\def\Picardfunctor{\mathrm{Pic}}
\def\Deformationcategory{\mathcal{D}\!\mathit{ef}}
\usepackage{cleveref}
\crefname{originalsection}{original section}{original sections}
\Crefname{originalsection}{Original section}{Original sections}
\makeatletter
\newlabel{derived-section-homotopy-triangulated}{{014P}{source}{}{sourcepointer-1}{}}
\newlabel{derived-section-homotopy-triangulated@cref}{{[originalsection][0][]014P}{[1][0][]source}{}{}{}}
\makeatother
\title{The sign in the boundary of a cup product of Cech cohomology classes}
\author{The Stacks Project Authors}
\date{stacks-project:cohomology.tex:local-snapshot-20261011; GNU FDL 1.2 or later}
\begin{document}
\maketitle
\noindent\textit{Editorial note.} All delivered mathematical text is the author's own, copied from the pinned Stacks Project chapter file cohomology.tex; no statement or proof step was rewritten, reordered or removed.\par\medskip
\noindent\textit{Editorial note.} The author attaches a footnote explaining the dependence of the sign on the conventions for the long exact sequence associated to a triangle in D(X); that footnote is part of the counted proof span and is reproduced as written.\par\medskip
\noindent\textit{Editorial note.} The counted spans carry 1 navigation stub and no bibliography citation.\par\medskip
\noindent\textit{Editorial note.} The preamble is the chapter's retained minimal preamble (sources/stacks-cohomology/retained-preamble.tex lines 1-81), a verbatim excerpt of the native preamble.\par\medskip
\noindent\textit{Editorial note.} The label lemma-compute-sign-cup-product-boundaries is defined only in cohomology.tex, so no cross-chapter restatement had to be weighed for this card.\par\medskip
\noindent\textit{Allowed standard prerequisites (editorial).}\begin{itemize}
\item Cech cohomology of an open covering with coefficients in a complex. Same chapter, the Cech complex of a covering with coefficients in a complex of sheaves, used throughout proof lines 5689-5715. Supplies the representatives of the classes a\_3 and b\_1 and the total complex in which the computation takes place.
\item Total complex of a double complex and the Leibniz rule for its differential. Homology, the sign conventions for the total complex, used at proof lines 5717-5723. Supplies the identity d(alpha cup beta) = d(alpha) cup beta + (-1)\textasciicircum{}n alpha cup d(beta).
\item Boundary map of a short exact sequence of complexes in cohomology. Same chapter, the boundary map attached to the short exact sequences of complexes, named in the statement at lines 5666-5668. Supplies the classes boundary(a\_3) and boundary(b\_1) that the computation identifies as the cocycles alpha\_1 and beta\_3.
\item Commutativity of the product with the maps of the short exact sequences. Same chapter, the commutative diagrams displayed immediately before the lemma in the source. Supplies gamma\_2(s tensor t) = gamma\_3(s-bar tensor t) and its companions, the substitution that produces the final two terms.
\end{itemize}\medskip
\ifdefined\hypertarget\hypertarget{sourcepointer-1}{}\fi
\noindent\textit{Source reference.} Original section 014P: sources/stacks-chow/tags.txt, source lines 1480--1480. Pointer outside the counted span, invoked once for context; the load-bearing argument of the counted statement is the author's proof reproduced in full in the proof span.\par\medskip
\begin{lemma}
\label{lemma-compute-sign-cup-product-boundaries}
In the situation above, assume {\v C}ech cohomology agrees with cohomology
for the sheaves $\mathcal{F}_i^p$ and $\mathcal{G}_j^q$.
Let $a_3 \in H^n(X, \mathcal{F}_3^\bullet)$ and
$b_1 \in H^m(X, \mathcal{G}_1^\bullet)$. Then we have
$$
\gamma_1( \partial a_3 \cup b_1) =
(-1)^{n + 1} \gamma_3( a_3 \cup \partial b_1)
$$
in $H^{n + m + 1}(X, \mathcal{H}^\bullet)$ where $\partial$ indicates the
boundary map on cohomology associated to the short exact sequences of
complexes above.
\end{lemma}

\begin{proof}
We will use the following conventions and notation. We think of
${\mathcal F}_1^p$ as a subsheaf of ${\mathcal F}_2^p$ and we think of
${\mathcal G}_3^q$ as a subsheaf of ${\mathcal G}_2^q$. Hence if $s$ is
a local section of ${\mathcal F}_1^p$ we use $s$ to denote
the corresponding section of ${\mathcal F}_2^p$ as well. Similarly
for local sections of ${\mathcal G}_3^q$. Furthermore,
if $s$ is a local section of ${\mathcal F}_2^p$ then we denote
$\bar s$ its image in ${\mathcal F}_3^p$. Similarly for the
map ${\mathcal G}_2^q \to {\mathcal G}^q_1$. In particular if
$s$ is a local section of ${\mathcal F}_2^p$ and $\bar s = 0$
then $s$ is a local section of ${\mathcal F}_1^p$. The commutativity
of the diagrams above implies, for local sections $s$ of
${\mathcal F}_2^p$ and $t$ of ${\mathcal G}_3^q$ that
$\gamma_2(s \otimes t) = \gamma_3(\bar s \otimes t)$ as sections of
${\mathcal H}^{p + q}$.

\medskip\noindent
Let ${\mathcal U} : X =  \bigcup_{i \in I} U_i$
be an open covering of $X$. Suppose that $\alpha_3$,
resp.\ $\beta_1$ is a degree $n$, resp.\ $m$ cocycle of
$\text{Tot}(
\check{\mathcal{C}}^\bullet({\mathcal U}, {\mathcal F}_3^\bullet))$,
resp.\ $\text{Tot}(
\check{\mathcal{C}}^\bullet({\mathcal U}, {\mathcal G}_1^\bullet))$
representing $a_3$, resp.\ $b_1$. After refining $\mathcal{U}$ if necessary,
we can find cochains $\alpha_2$, resp.\ $\beta_2$ of
degree $n$, resp.\ $m$ in
$\text{Tot}(
\check{\mathcal{C}}^\bullet({\mathcal U}, {\mathcal F}_2^\bullet))$,
resp.\ $\text{Tot}(
\check{\mathcal{C}}^\bullet({\mathcal U}, {\mathcal G}_2^\bullet))$
mapping to $\alpha_3$, resp.\ $\beta_1$.
Then we see that
$$
\overline{d(\alpha_2)} = d(\bar \alpha_2) = 0
\quad\text{and}\quad
\overline{d(\beta_2)} = d(\bar \beta_2) = 0.
$$
This means that $\alpha_1 = d(\alpha_2)$ is a degree $n + 1$ cocycle in
$\text{Tot}(\check{\mathcal{C}}^\bullet({\mathcal U}, {\mathcal F}_1^\bullet))$
representing $\partial a_3$. Similarly, $\beta_3 = d(\beta_2)$ is
a degree $m + 1$ cocycle in
$\text{Tot}(\check{\mathcal{C}}^\bullet({\mathcal U}, {\mathcal G}_3^\bullet))$
representing $\partial b_1$.
Thus we may compute
\begin{align*}
d(\gamma_2(\alpha_2 \cup \beta_2))
& =
\gamma_2(d(\alpha_2 \cup \beta_2))
\\
& =
\gamma_2(d(\alpha_2) \cup \beta_2 + (-1)^n \alpha_2 \cup d(\beta_2) )
\\
& =
\gamma_2( \alpha_1 \cup \beta_2)  + (-1)^n \gamma_2( \alpha_2 \cup \beta_3)
\\
& =
\gamma_1(\alpha_1 \cup \beta_1) + (-1)^n \gamma_3(\alpha_3 \cup \beta_3)
\end{align*}
So this even tells us that the sign is $(-1)^{n + 1}$ as indicated
in the lemma\footnote{The sign depends on the convention for the
signs in the long exact sequence in cohomology associated to a triangle
in $D(X)$. The conventions in the Stacks project are (a) distinguished
triangles correspond to termwise split exact sequences and (b) the boundary
maps in the long exact sequence are given by the maps in the snake lemma
without the intervention of signs. See
Derived Categories, Section \ref{derived-section-homotopy-triangulated}.}.
\end{proof}

\end{document}
